\documentclass[10pt,a4paper]{amsart}
\usepackage[margin=19mm]{geometry}
\usepackage{amsmath,amssymb,mathtools}
\usepackage[scr=rsfs,cal=euler]{mathalfa}
\usepackage{xfrac}
\usepackage[unicode,hidelinks,bookmarks=false]{hyperref}
\newcommand{\ie}{i.e.\ }
\newcommand{\cf}{cf.\ }
\newcommand{\resp}{respectively\ }
\newcommand{\Subsection}{\subsection}
\providecommand{\nobreakdash}{\nobreak\hbox{-}}
\setlength{\emergencystretch}{2em}
\multlinegap0pt
\usepackage{amssymb}
\usepackage[shortlabels]{enumitem}
\usepackage{float}
\usepackage{tikz}
\newtheorem{lemma}{Lemma}
\title{Domination and packing in graphs}
\author{\'Akos D\'ucz, Anna Gujgiczer}
\date{Source: arXiv:2602.18402v2 (CC BY 4.0)}
\begin{document}
\maketitle

\noindent\emph{Source and licence.} Domination and packing in graphs. Version arXiv:2602.18402v2 (CC BY 4.0), distributed by arXiv under the Creative Commons Attribution 4.0 International licence, http://creativecommons.org/licenses/by/4.0/, as recorded in the arXiv licence metadata for this exact version and shown on the abstract page https://arxiv.org/abs/2602.18402v2. Full licence notice: \texttt{licenses/arxiv-2602.18402-CC-BY-4.0.txt}.\par\smallskip

\noindent\textit{Editorial scope.} Counted unit: the article's lemma lem:discharge (source0.tex lines 321--323). The article's complete original proof of it is delivered with it (source0.tex lines 325--333, one \texttt{proof} environment of 9 lines). No mathematics is written, repaired or re-derived: the statement and the proof are the article's own lines, in the article's own notation, and no retained line is substituted or re-flowed.  Retained uncounted as the source-local context the proof needs: lines 307--314.  Nothing was omitted from any retained span, so the package carries no omission record.  No retained line contains a citation, so the wrapper carries no bibliography and no bibliography entry.  No retained line contains a figure environment, an image-inclusion command or drawing code, so no image is delivered and the package declares no asset dependency.  Editorial formatting only, in addition: an AMS \texttt{amsart} preamble that restates the article's own declarations and macros used by the retained lines, namely the environments \texttt{lemma} on separate counters, together with the standard packages those lines need.  The retained environments print in this document as Lemma 1 in order of appearance, whereas the article prints the same items with the numbering of its own full document.  The complete native source of the article as distributed in the arXiv e-print for this exact version is bundled unchanged as \texttt{sources/195037/source0.tex}.
\begin{lemma}\label{lem:triangulate}
    If $H$ is a simple connected planar graph and $I$ is and independent set in $H$, then we can add new edges to $H$ such that:
    \begin{itemize}
        \item $I$ is still independent in the new graph $H'$.
        \item $H'$ is planar (but not necessarily simple).
        \item Every face of $H'$ is a triangle.
    \end{itemize}
\end{lemma}

\begin{lemma}\label{lem:discharge}
    Let $H$ be a simple planar graph with minimum degree $4$. Then $H$ has an edge $(x, y)$ with $d(x) \leq 7$ and $d(y) \leq 7$.
\end{lemma}

\begin{proof}
    We may assume that $H$ is connected (otherwise we just take some component of $H$).
    Let $I$ denote the set of vertices in $H$ with $d(v) \leq 7$. If $I$ is not independent, the statement follows. Suppose that $I$ is independent. By Lemma \ref{lem:triangulate}, we can add new edges to $H$ to obtain $H'$, every face of which is a triangle and in which $I$ is still independent.

    Let us assign a charge of $d(v) - 6$ to every vertex $v$ of $H'$. By Euler's formula, we have that the total sum of these charges is negative. 
    Now let each vertex $v$ with $d(v) \geq 8$ give a charge of $1/2$ to each of it's neighbors in $I$ (if there are multiple edges between $v$ and $w$, then we count $w$ as a 'neighbor' multiple times). Since $H'$ is triangulated and $I$ is independent in $H'$, such a vertex $v$ gives charge to at most half of it's neighbors. Thus every vertex with $d(v) \geq 8$ has nonnegative charge after redistribution.

    If $d(v) \leq 7$ then all of its neighbors have degree at least $8$, therefore $v$ must receive $1/2$ charge from each neighbor. The initial charge of $v$ is smallest in case $d(v) = 4$, which gives a charge of $d(v)-6 + 4\cdot(1/2) = 0$ after discharging. Therefore every vertex has nonnegative charge at the end of the procedure, which is a contradiction. 
\end{proof}

\end{document}
